\documentclass[12pt,a4paper]{article}
\usepackage[utf8]{inputenc}
\usepackage[T1]{fontenc}
\usepackage[ngerman]{babel}
\usepackage{amsmath,amssymb}
\usepackage{geometry}
\geometry{a4paper,left=2.5cm,right=2.5cm,top=2.5cm,bottom=2.5cm}
\usepackage{parskip}

\title{\textbf{Warum die Allgemeine Relativitätstheorie prinzipiell unvollständig sein muss} \\
\large — und was an ihre Stelle tritt}
\author{Ein Diskussionsbeitrag}
\date{}

\begin{document}

\maketitle

\section*{Vorbemerkung}

Dieser Text ist kein Angriff auf die Erfolge der Allgemeinen Relativitätstheorie. Sie hat sich in unzähligen Experimenten bewährt und ist eine der präzisesten Theorien der Physik. Aber Präzision allein garantiert keine prinzipielle Richtigkeit. Die Frage ist: Ruht die ART auf einem Fundament, das wirklich trägt? Oder gibt es einen tiefer liegenden Grund, warum sie funktioniert – und gleichzeitig einen Punkt, an dem sie prinzipiell unvollständig sein muss?

Die folgende Argumentation zeigt: Die ART ist nicht die fundamentale Theorie der Gravitation. Sie ist ein emergenter Grenzfall einer einfacheren, tiefer liegenden Theorie – der Dynamischen Schwere-Trägheits-Theorie (DSTT).

\section{Das Äquivalenzprinzip: Ein Fundament mit Rissen}

Die Allgemeine Relativitätstheorie ruht auf einem Fundament: dem \textbf{Äquivalenzprinzip}. Es besagt, dass ein homogenes Gravitationsfeld lokal nicht von einem beschleunigten Bezugssystem unterscheidbar ist. Oder anders formuliert: Träge Masse und schwere Masse sind identisch.

Dieses Prinzip ist bestechend einfach. Aber es enthält eine konzeptuelle Schwäche, die selten hinterfragt wird: \textbf{Es vermischt Ursache und Wirkung.}

\begin{itemize}
    \item Die \textbf{Ursache} ist die Gravitation – eine äußere Einwirkung auf einen Körper.
    \item Die \textbf{Wirkung} ist die Trägheit – die Art und Weise, wie der Körper darauf reagiert.
\end{itemize}

Die ART setzt beides gleich. Sie sagt: Die Reaktion eines Körpers auf Gravitation ist identisch mit seiner Reaktion auf Beschleunigung. Das ist eine \textbf{ungeprüfte Identifikation} – und sie ist, wie sich zeigen lässt, \textbf{nicht haltbar}.

\section{Die DSTT: Trennung von Ursache und Wirkung}

Die \textbf{Dynamische Schwere-Trägheits-Theorie (DSTT)} stellt die Dinge richtig. Ihre Grundaxiome sind denkbar einfach:

\begin{enumerate}
    \item \textbf{Schwere ist die einzige Ursache.} Sie wirkt radial zwischen zwei Körpern. Es gibt keine zusätzlichen Felder, keine Raumzeitkrümmung – nur diese eine radiale Wechselwirkung.
    
    \item \textbf{Trägheit ist die einzige Wirkung.} Sie zerfällt in zwei Komponenten: eine radiale und eine tangentiale Antwort. Die Schwerewirkung teilt sich auf in einen radialen Anteil \((1-\beta)F_G\) und einen tangentialen Anteil \(\beta F_G\).
    
    \item \textbf{Die Aufteilung ist dynamisch.} Ein Zustandsparameter \(\beta(t) \in [0,1]\) bestimmt das Verhältnis. Er ist definiert als das Verhältnis von tangentialer zu Gesamtenergie:
    \[
    \beta = \frac{E_B}{E} = \frac{\frac{1}{2}m_T r^2\dot{\phi}^2}{\frac{1}{2}m_T\dot{r}^2 + \frac{1}{2}m_T r^2\dot{\phi}^2}
    \]
    
    \item \textbf{\(\beta(t)\) wächst monoton.} Aus den Energieflüssen folgt zwingend:
    \[
    \dot{\beta} > 0
    \]
    Die tangentiale Komponente nimmt ständig zu, erreicht aber nie den Wert 1.
\end{enumerate}

Mehr braucht es nicht. Keine gekrümmte Raumzeit, keine dunkle Materie, keine zusätzlichen Felder – nur diese vier Aussagen.

\section{Die unmittelbaren Konsequenzen}

\subsection{Es gibt keine geschlossenen Bahnen}

Aus den Axiomen folgt durch Eliminierung der Zeit:
\[
\frac{dr}{d\phi} = \frac{r}{B}, \quad \text{mit } B = \frac{2E}{F_G}
\]

Die Lösung ist eine \textbf{logarithmische Spirale}:
\[
r(\phi) = K e^{\phi/B}
\]

Kreisbahnen und Ellipsen – die Grundbausteine der newtonschen Mechanik und der ART – existieren nicht. Sie sind nur Näherungen für kurze Zeiträume, in denen die Drift vernachlässigbar erscheint.

\subsection{Gravitation ist irreversibel}

Weil \(\dot{\beta} > 0\) und damit \(\dot{r} \neq 0\), ist jede Bewegung mit einer \textbf{Drift} verbunden. Es gibt keine zeitumkehrsymmetrischen Bahnen. Gravitation produziert permanent \textbf{Entropie}. Sie definiert einen fundamentalen \textbf{Zeitpfeil} – nicht abgeleitet aus der Thermodynamik, sondern direkt aus der Dynamik der Gravitation selbst.

\subsection{Langfristig driftet alles nach außen}

Da \(\beta(t)\) monoton wächst, wird jede anfängliche Einwärtsbewegung (\(\dot{r} < 0\)) irgendwann in eine Auswärtsbewegung (\(\dot{r} > 0\)) umgekehrt. Der Umschlagspunkt ist durch die Anfangsbedingungen bestimmt, aber der langfristige Trend ist eindeutig:

\textbf{Alle Satelliten, Planeten, Monde – alles spiralisiert langfristig nach außen.}

\section{Was das für die ART bedeutet}

Die ART basiert auf drei Säulen:

\begin{enumerate}
    \item \textbf{Das Äquivalenzprinzip} (Ursache = Wirkung)
    \item \textbf{Die gekrümmte Raumzeit} (Geometrie als Ursache der Bewegung)
    \item \textbf{Die Feldgleichungen} (Materie krümmt Raumzeit, Raumzeit bestimmt Bewegung)
\end{enumerate}

Die DSTT zeigt: \textbf{Säule 1 ist falsch.} Ursache und Wirkung sind nicht identisch. Die Trägheitsreaktion ist anisotrop – sie zerfällt in zwei Komponenten. Das Äquivalenzprinzip ist eine Idealisierung, die die wahre Struktur der Gravitation verschleiert.

Wenn aber das Fundament der ART falsch ist, dann ist die Theorie als Ganze \textbf{prinzipiell unvollständig}. Sie kann nicht die fundamentale Theorie der Gravitation sein – allenfalls eine Näherung, ein \textbf{emergenter Grenzfall}.

\section{Die ART als emergenter Grenzfall}

Wann funktioniert die ART? Wenn:

\begin{itemize}
    \item \(\beta \approx \text{konstant}\) (keine merkliche Drift über kurze Zeiträume)
    \item \(\dot{r} \approx 0\) (quasi-geschlossene Bahnen)
    \item Die anisotrope Trägheit vernachlässigt werden kann
\end{itemize}

In diesem Grenzfall kann man die Spiralbahn durch eine präzedierende Ellipse annähern. Die ART beschreibt diese Präzession – aber sie übersieht die zugrundeliegende Drift. Sie ist eine \textbf{Statik, die eine Dynamik vortäuscht}.

Vergleichbar ist das mit der Newtonschen Mechanik: Sie ist ein Grenzfall der ART für kleine Geschwindigkeiten und schwache Felder. Niemand würde behaupten, Newton sei "falsch" – aber er ist unvollständig. Genauso ist die ART unvollständig: Sie ist der Grenzfall der DSTT für vernachlässigbare \(\dot{\beta}\).

\section{Was die DSTT ohne Zusatzannahmen erklärt}

Die DSTT allein – ohne Dunkle Materie, ohne gekrümmte Raumzeit, ohne zusätzliche Felder – erklärt eine Reihe von Phänomenen, die in der Standardphysik als Probleme gelten:

\begin{tabular}{|l|l|l|}
\hline
\textbf{Phänomen} & \textbf{Bisherige Erklärung} & \textbf{DSTT-Erklärung} \\
\hline
Merkur-Periheldrehung & ART (43'' pro Jahrhundert) & Natürliche Folge der Spiralbahn \\
Jupiterwanderung & Komplexe Resonanzen & Universelle Auswärtsdrift \\
Hot Jupiters & Einwanderung durch Reibung & Frühe Einwärtsphase \\
Galaxienrotation & Dunkle Materie & Anisotrope Trägheit + \(\beta\)-Wachstum \\
Monddrift & Gezeitenreibung (zu klein) & Fundamentale Drift durch DSTT \\
Phobos & Sollte einwärts fallen & Einwärtsphase, langfristig umkehrbar \\
\hline
\end{tabular}

Keine dieser Erklärungen benötigt zusätzliche Hypothesen. Sie folgen direkt aus den vier Axiomen der DSTT.

\section{Die entscheidende Frage}

Warum wurde das nicht längst erkannt? Weil die Drift extrem langsam ist. Für die Erde um die Sonne beträgt sie:
\[
\frac{\dot{r}}{r} \approx 10^{-11} \text{ pro Jahr}
\]

Das ist in historischen Zeiträumen nicht messbar. Ein Erdumlauf dauert ein Jahr – die Drift pro Umlauf ist etwa \(10^{-11}\) des Bahnradius. Das entspricht wenigen Millimetern pro Jahr. Kein Wunder, dass Kepler und Newton geschlossene Ellipsen sahen.

Aber auf kosmischen Skalen summiert sich diese Drift. Und genau das sehen wir:

\begin{itemize}
    \item Jupiter ist nach außen gewandert – um ein Vielfaches mehr, als Gezeitenmodelle erklären können.
    \item Der Mond driftet nach außen – schneller als Gezeitenmodelle vorhersagen.
    \item Galaxien zeigen flache Rotationskurven – als hätten sie zusätzliche Masse, aber sie haben nur eine andere Trägheitsdynamik.
    \item Das Universum zeigt eine Rotverschiebung, die wie Expansion aussieht – aber wenn alle Bahnen Spiralen sind, dann addieren sich die Driften zu einer scheinbaren Expansion.
\end{itemize}

\section{Der entscheidende Test}

Die DSTT macht eine klare, von der ART abweichende Vorhersage: \textbf{Alle Bahnen sind Spiralen.} Das bedeutet:

\begin{enumerate}
    \item Die Drift des Mondes ist \textbf{nicht primär durch Gezeiten} verursacht, sondern durch die Gravitation selbst. Präzise Lasermessungen sollten eine systematische Komponente zeigen, die sich nicht durch Gezeitenmodelle erklären lässt.
    
    \item Die Jupiterwanderung ist \textbf{kein Zufall} der Planetenentstehung, sondern ein universelles Phänomen. Auch andere Gasriesen sollten systematische Wanderung zeigen.
    
    \item Präzise Satellitenbahnen (wie die der Mondsonde LRO oder der Mars-Sonden) müssen eine minimale, aber systematische radiale Drift zeigen – wenn man lange genug misst.
    
    \item Die Rotverschiebung ist \textbf{frequenzabhängig}. Da \(\beta\) für Photonen von der Energie abhängt, müssen verschiedene Frequenzen leicht unterschiedliche Rotverschiebungen zeigen. Dieser Effekt ist klein, aber mit Pulsar-Laufzeiten oder hochpräzisen Spektren messbar.
\end{enumerate}

\section{Fazit}

Die ART ist nicht falsch im Sinne von "liefert falsche Zahlen". Für viele Zwecke ist sie eine hervorragende Näherung. Aber sie ist \textbf{prinzipiell unvollständig}, weil sie auf einer falschen Identifikation beruht: der Gleichsetzung von Ursache und Wirkung.

Die DSTT stellt die Dinge richtig:

\begin{itemize}
    \item \textbf{Ursache} = Schwere (radiale Wirkung)
    \item \textbf{Wirkung} = Anisotrope Trägheit (radial + tangential)
    \item \textbf{Dynamik} = Monoton wachsender Umlenkungszustand \(\beta(t)\)
    \item \textbf{Bahnform} = Logarithmische Spirale
\end{itemize}

Das ist keine komplizierte Theorie. Es sind \textbf{vier Axiome}, die sich auf einem Blatt Papier notieren lassen. Und sie erklären mehr als die ART – ohne Dunkle Materie, ohne gekrümmte Raumzeit, ohne Singularitäten, ohne Renormierung.

Die ART emergiert aus der DSTT, wenn man die Drift vernachlässigt. Aber das Fundament ist tiefer. Und dieses Fundament heißt:

\begin{center}
\Large
\textbf{Ursache und Wirkung sind nicht dasselbe.}
\end{center}

\section*{Ausblick}

Wenn dieses Prinzip stimmt, dann ist die ART nicht das Ende, sondern ein Anfang. Sie ist der sichtbare Teil eines viel größeren Gebäudes – eines Gebäudes, das mit der DSTT beginnt und über die Weber-Gravitation und das Quantenpotential bis hin zu einer vollständigen, deterministischen, singularitätsfreien Physik reicht.

Die DSTT ist der einfachste Zugang zu diesem Gebäude. Sie zeigt: Die Gravitation ist keine Geometrie. Sie ist Dynamik. Und diese Dynamik ist einfacher, als wir dachten.

\end{document}